\chapter{Base de conhecimento}
\label{cap:base-conhecimento}

A base de conhecimento (\cod{driving_rules.pl}) é um conjunto de $\sim$120 regras de
produção sobre condução, que replicam as regras de decisão do motor Drools do
sistema.

\section{O que a base contém}
\label{sec:base-conteudo}

O objetivo é que o comportamento do veículo fique descrito nas regras, e não
espalhado pelo código de controlo. A base cobre emergências, peões, veículo da
frente e ultrapassagens, limites de velocidade, STOP, semáforos, perigos na estrada,
mudanças de via, barreiras e meteorologia. Várias regras seguem diretamente o Código da Estrada~\cite{codigoestrada}: o
Art.~13.º, n.º~3 (circular pela via mais à direita, \cref{lst:regra-direita}),
o Art.~18.º (distância ao veículo da frente), o Art.~24.º (adequar a velocidade
às condições) e os Arts.~36.º e 38.º (ultrapassar pela esquerda e regressar à
direita no fim da manobra).

As regras leem $\sim$66 tipos de facto, que se agrupam como mostra a
\cref{tab:tipos-factos}.

\begin{table}[H]
  \centering
  \caption{Tipos de facto lidos pelas regras}
  \label{tab:tipos-factos}
  \small
  \begin{tabularx}{\textwidth}{@{}L{3.0cm}V>{\raggedright\arraybackslash}XVL{4.2cm}@{}}
    \toprule
    \textbf{Tipo} & \textbf{Exemplos} & \textbf{Origem} \\
    \midrule
    Perceção & \cod{ego_speed(Kmh)}, \cod{pothole(Id, DistanceM, Severity, Lane)},
        \cod{traffic_light(Id, red|yellow|green, DistanceM)},
        \cod{lane_occupancy(Side, Occupied, FrontDistanceM)}
      & Sensores e veículo, carregados por \cod{driving_state.pl} no início do
        ciclo \\ \linhasep
    Estado anterior & \cod{overtaking/5}, \cod{light_commitment/2},
        \cod{stop_hold_started/2}, \cod{previous_decision/2}
      & Estado guardado entre ciclos (38 tipos) e a decisão anterior \\ \linhasep
    Conclusões intermédias & \cod{pedestrian_conflict(Id)},
        \cod{fuzzy_membership/3}, \cod{fuzzy_conclusion/2},
        \cod{overtake_started/3}
      & Concluídas por regras no próprio ciclo \\ \linhasep
    Propostas & \cod{proposal(Action, SpeedKmh, Risk, Priority, Reason)}
      & Concluídas por regras, arbitradas no Java \\
    \bottomrule
  \end{tabularx}
\end{table}

Além das regras e dos factos, a base tem predicados auxiliares, que leem os
factos através de \cod{known/1} (por exemplo \cod{front_vehicle_within/1}), e
subobjetivos provados por \emph{backward chaining} (\cref{sec:base-exemplos}). As
constantes numéricas, como as distâncias de paragem e o valor de cada
prioridade, não estão escritas no Prolog: chegam do Java no pedido
\cod{reset} (\cod{param/2}, \cod{priority_value/2}).

Dada a extensão recomendada para este relatório (cerca de 10 páginas), não se
inclui o diagrama completo da árvore de decisão, isto é, o fluxograma de decisão
da velocidade e da via.
Esse diagrama pode
ser consultado junto ao código-fonte do projeto (\cref{anx:mapeamento}).
O percurso que cada decisão segue nesse fluxograma pode também ser consultado
na \emph{interface} Web, através da explicação \emph{How} (\cref{sec:expl-how}).

\section{Representação e organização}
\label{sec:base-representacao}

O ficheiro \cod{driving_rules.pl} está organizado por situação (emergência,
peões, semáforos, sinal STOP, veículo da frente, ultrapassagem, perigos na
estrada, barreiras, meteorologia, velocidade, arranque e decisão base), seguido
dos subobjetivos e dos predicados auxiliares. A precedência acompanha a
urgência: as regras de emergência têm 870 a 900, as de peões 600 a 1000, as de
velocidade 205 a 310, e a decisão base é a última (60 a $-100$). Nos termos da
unidade curricular, a base tem regras de ação, que concluem uma \cod{proposal/5} (a
maioria), e regras dedutivas, que concluem factos intermédios como
\cod{pedestrian_conflict/1} ou \cod{fuzzy_conclusion/2}; a negação entra nas
condições, com \cod{not/1}, e não há regras de refutação: um facto nunca é
negado, apenas ausente. A \cref{lst:regra-limite-pl} mostra uma regra de ação.

\begin{lstlisting}[caption={Regra de excesso de velocidade em Prolog (\texttt{driving\_rules.pl})},label=lst:regra-limite-pl]
rule('Reduce speed when exceeding detected limit',
     300,
     [ ego_speed(V),
       speed_limit(Limit),
       test(V > Limit + 0.1) ],
     [ assert(proposal(reduce_speed, Limit, moderate, speed_compliance,
                       'Vehicle exceeds current legal speed limit')) ]).
\end{lstlisting}

Ou seja, se a velocidade passar o limite em vigor em mais de 0,1\,km/h,
propõe-se reduzir para o limite, com risco moderado e prioridade de cumprimento
de velocidade (20).


\section{Alguns exemplos}
\label{sec:base-exemplos}

\paragraph{Circular pela direita.} A \cref{lst:regra-direita}, uma das regras mais longas da base, aplica o Art.~13.º, n.º~3,
do Código da Estrada (circular pela via mais à direita) e usa todas as formas de
condição da \cref{sec:motor-regras}. Só propõe voltar à direita quando nada mais
exige atenção: nenhum obstáculo, buraco, barreira ou
peão em conflito, nenhum veículo a menos de 80\,m e nenhuma ultrapassagem ou
mudança de via pedida pelo condutor.

\noindent\begin{minipage}{\linewidth}
\begin{lstlisting}[caption={Uma das regras mais longas da base de conhecimento (\texttt{driving\_rules.pl})},label=lst:regra-direita]
rule('Enforce keep right road law when right lane is safe',
     460,
     [ ego_lane(Lane),                           test(Lane \== right),
       not(presentation_demo),                   lane_change_state(idle),
       prove(no_continuous_line(right)),         not(right_lane_blocked_ahead),
       time_ms(Now),                             test(Now >= 4000),
       lane_occupancy(right, false, RightFront), test(RightFront > 80.0),
       test(fresh_side_reading(right)),          not(move_off_state(true, _)),
       not(obstacle(_, _, _, _)),                not(pothole(_, _, _, _)),
       not(barrier(_, _, false, _)),             not(pedestrian_conflict(_)),
       not(test(front_vehicle_within(80.0))),    not(lane_change_request(_)),
       not(test(driver_lane_change_hold(Now))),
       not(test(current_overtaking(true, _, _, _, _))) ],
     [ assert(proposal(return_right, no_target, low, maneuver_safety,
       'Portuguese Highway Code Art. 13: Drive on the rightmost available safe lane')) ]).
\end{lstlisting}
\end{minipage}

\paragraph{Limite de velocidade.} A câmara deteta o sinal até 120\,m à frente.
O \cod{driving_state.pl} guarda-o como pendente, com a sua posição na estrada, e
carrega-o como \cod{pending_speed_limit(Kmh, SignPositionM)}. Enquanto o sinal
está à frente, \cod{'Slow down in time for the speed limit ahead'} faz o veículo
chegar a ele já ao novo limite. Quando o veículo chega ao sinal, o limite
pendente passa a \cod{speed_limit(Kmh)}, e a regra da
\cref{lst:regra-limite-pl} pode propor \cod{reduce_speed}. Os dois limites estão
separados porque, aplicado logo que o sinal era visto, o limite fazia o veículo
abrandar demasiado cedo.

\paragraph{Meteorologia.} Para a meteorologia recorre-se a lógica difusa: um
controlador de Mamdani com 8 regras, em três níveis de precedência. Duas regras
de fuzificação (360) convertem a visibilidade e a aderência de
\cod{weather/3} em graus de pertença, \cod{fuzzy_membership(Variable, Term, Degree)}. Cinco regras de inferência (355) combinam-nos, com E como mínimo e OU
como máximo, e concluem \cod{fuzzy_conclusion(Term, Degree)}
(\cref{lst:regra-difusa}). Uma regra de desfuzificação (350) agrega as
conclusões e calcula o centroide, que é a velocidade recomendada: um único teto
\cod{reduce_speed} com prioridade de perigo na estrada (40). A precedência garante
que cada nível só dispara depois do anterior.

\begin{lstlisting}[caption={Uma regra de inferência difusa (\texttt{driving\_rules.pl})},label=lst:regra-difusa]
rule('Fuzzy: poor visibility or poor grip -> very slow',
     355,
     [ fuzzy_membership(visibility, poor, V),
       fuzzy_membership(grip, poor, G),
       test(( D is max(V, G), D > 0.0 )) ],
     [ assert(fuzzy_conclusion(very_slow, D)) ]).
\end{lstlisting}

\paragraph{Mudança de via segura.} \cod{safe_to_change_lane/1} é um objetivo,
provado por \emph{backward chaining} a partir de quatro subobjetivos
(\cref{lst:safe-lane}). As regras de desvio chamam-no nas condições com
\cod{prove(safe_to_change_lane(left))}. Mudar de via é seguro se não houver
outra manobra em curso nem linha contínua, se o ângulo morto estiver livre e se
houver espaço para seguir o veículo dessa via, com 12\,m de folga
(\cod{StoppingPolicy.roomToFollow}).

\begin{lstlisting}[caption={Objetivo provado por \emph{backward chaining} (\texttt{driving\_rules.pl})},label=lst:safe-lane]
safe_to_change_lane(Side) :-
    maneuver_idle,
    no_continuous_line(Side),
    side_clear(Side),
    room_to_follow(Side).
\end{lstlisting}

\paragraph{Incerteza.} A perceção é tratada como incerta de três formas: as
leituras vagas, como a visibilidade e a aderência, passam pelo controlador
difuso; as leituras ruidosas (o radar frontal tem ruído gaussiano com desvio
padrão de 0,7\,m) só confirmam uma situação crítica quando se repetem, por
exemplo três amostras de TTC crítico numa janela de 2\,s na regra de travagem
de emergência; e as leituras que deixam de chegar são esquecidas
(\cref{sec:base-manutencao}).

\section{Acesso e manutenção}
\label{sec:base-manutencao}

Não há manutenção de verdade: cada ciclo recarrega os factos a partir do estado em
\cod{driving_state.pl}, que esquece o que deixou de ser observado
(\cod{expire_old_readings/1}). As regras só mudam editando os ficheiros.
